%-------------------------------------------------------------------------------
%	SECTION TITLE
%-------------------------------------------------------------------------------
% 在 main.tex 或 cv/experience.tex 中使用

% --- 科研经历开始 ---
% ================= 科研经历 START =================


% ====== 项目 1：Type-I (一作) ======
 \cvsection{科研经历}

\begin{cventries}

% ====== 统一大项目：天津市自然科学基金 ======
 \cventry
    {项目主要参与者（核心研究员） $|$ 负责人(导师)：王景 副教授} 
    {天津市自然科学基金面上项目：两维二次型能带交叉费米体系中四费米子作用及杂质散射的低能物理效应} 
    {\vspace{1em}项目编号：25JCYBJC01640} 
    {} % 基金行不再写时间
    {
      \begin{cvitems}
        % --- 课题 1 ---
        \vspace{0.5mm}
        \item[] \hspace*{-1em}{\textbf{课题一：二维 QBCP 系统中的 2+1 维 QED 有效场论与重整化群研究（第一作者）} \hfill 2024.09 -- 至今}
          \begin{itemize}
            \item {\textbf{研究背景与理论建模}：旨在探究纯费米子体系引入规范场（光子）后的新奇物理效应，\textbf{将 3+1 维 QED 理论框架推广至 2+1 维凝聚态 QBCP 体系}。针对基于棋盘晶格（Checkerboard lattice）的 $z=2$ 二次能带交叉点半金属，针对棋盘晶格半金属构建了包含费米子-光子耦合与多通道四费米子短程相互作用的有效场论，通过重整化群（RG）方法系统分析了体系的涌现行为。}
            \item {\textbf{单圈计算与构建RG 流方程}：独立完成体系在单圈修正层级的19个费曼图的解析积分与动量壳层微扰计算，基于场论微扰计算的结果推导建立包含 10 个非线性参数演化的重整化群（RG）流方程组，在理论上自洽统一了规范涨落与电子关联效应。}
            \item {\textbf{数值模拟与相图分析}：编写 MATLAB（ode15s 代码求解高维流方程；重点考察了没有QED时四费米子相互作用参数的固定点、QED参数的固定点、QED 效应对固定点漂移的影响。运用 \textbf{Origin} 绘制相图，展现固定点。}
            \item {\textbf{学术产出}：运用 LaTeX 排版逾 \textbf{70 页}的高密度解析推导与数据分析笔记；正进行论文初稿准备，\textbf{拟投递 Physical Review B}。}
          \end{itemize}
        \vspace{0.8em}
% --- 课题 2 ---
        \item[] \hspace*{-1em}{\textbf{课题二：第二类（Type-II）QBCP 半金属中的 QED 相互作用效应研究（第一作者）} \hfill 2025.11 -- 至今}
          \begin{itemize}
            \item {\textbf{模型推广}：将理论前沿拓展至笼目晶格（Kagome lattice），引入粒子--空穴对称性破缺项 $t_{I}k^{2}\Sigma_{1}$，将 2+1 维 QED 理论成功拓展至 Type-II QBCP 体系。}
            \item {\textbf{核心工作}：主导了费米子传播子与 QED 顶角函数的推导，负责计算各向异性参数在倾斜能带结构下的重整化群流向。}
            \item {\textbf{当前进度}：已完成有效理论构建，正在进行解析计算。}
          \end{itemize}
      \end{cvitems}
    }

\vspace{1em}

% --- 项目 3：中科院物理所实习 ---
  \cventry
    {科研实习生 $|$ 导师：鲁年鹏 研究员（SF03组）} 
    {中国科学院物理研究所（IOP CAS） $|$ 表面物理国家重点实验室} 
    {北京} 
    {2025.07 -- 2025.09} 
    {
      \begin{cvitems}
        \item {课题： 镍基超导薄膜的生长与调控。}
        \item {\textbf{薄膜生长}：操作 PLD 系统在 800°C 下通过系统调节氧压（8--16 Pa）与激光能量（36--46 mJ）优化 $La_2SmNi_2O_7$ 沉积参数 ，随后降温至 500°C 原位制备$SrTiO_3$ 盖帽层 ；样品经 XRD 鉴定后 进行退火处理，最后通过低温输运测量分析其电阻-温度转变特性。}
        \item {\textbf{表征与器件}：运用 pld长膜、XRD 进行物相鉴定；观摩学习无掩膜光刻技术与输运测量以及AFM的使用 。}
      \end{cvitems}
    }

\end{cventries}
% ================= 科研经历 END =================
% --- 科研经历结束 ---
